\begin{table}[H]
  \centering
  \caption{Matrices de transición de pobreza multidimensional (IPM)
  (\% de fila, panel emparejado 1 a 1)}
  \label{tab:matriz_transicion_ipm}
  \begin{tabular}{llcc}
    \toprule
    \textbf{Periodo} & \textbf{Estado inicial} & \textbf{No pobre (final)} & \textbf{Pobre (final)} \\
    \midrule
    \multirow{2}{*}{2010 $\rightarrow$ 2013 ($n=8{,}218$)}
      & No pobre & 87.9\% & 12.1\% \\
      & Pobre    & 54.5\% & 45.5\% \\
    \addlinespace
    \multirow{2}{*}{2013 $\rightarrow$ 2016 ($n=6{,}911$)}
      & No pobre & 90.3\% & 9.7\% \\
      & Pobre    & 52.3\% & 47.7\% \\
    \bottomrule
  \end{tabular}
  \begin{minipage}{0.85\textwidth}
    \vspace{4pt}
    \footnotesize \textit{Nota:} $n$ = hogares con emparejamiento 1 a 1;
    excluye 511 (2010--2013) y 1{,}190 (2013--2016) casos de hogares que se
    dividieron entre olas -- mismo criterio de emparejamiento que la matriz
    de pobreza monetaria (Tabla~\ref{tab:matriz_transicion}), aquí aplicado
    a la clasificación IPM. Fuente: cálculos propios con
    base en ELCA 2010, 2013 y 2016.
  \end{minipage}
\end{table}